\section{Limitations}

Despite the promising results, our work has several limitations that open avenues for future research.

\noindent \textbf{Architecture Dependence.}
Our experiments are conducted primarily on a YOLO-based detector (YOLOv10). While the proposed approach is designed to be model-agnostic, we do not validate its effectiveness across other detection architectures such as Faster R-CNN or transformer-based models. This limitation is partly due to computational constraints during experimentation. Evaluating the proposed framework across diverse architectures like Faster R-CNN remains an important direction for future work.

\noindent \textbf{Approximate Nature of Explainability Methods.}
The explainability techniques employed in this study, including Grad-CAM and occlusion sensitivity, provide approximate interpretations of model behavior. These methods do not guarantee causal attribution and may not fully capture the true decision-making process of the model. As a result, the insights derived from XAI analyses should be interpreted with caution.

\noindent \textbf{Limited Dataset Diversity.}
Our evaluation is restricted to specific in-distribution and out-of-distribution datasets (Pascal VOC and Caltech-256). However, real-world OoD scenarios are significantly more diverse and harder to systematically capture. The lack of standardized and representative OoD benchmarks limits the generalizability of our findings to broader real-world settings.

\noindent \textbf{Lack of Real-time Deployment Analysis.}
Although the proposed method is lightweight and designed for inference-time application, we do not perform a detailed evaluation on real-time systems or streaming data. In particular, the absence of deployment on live or production-level pipelines limits our understanding of its practical performance in real-world applications.